A cell's measured transcriptome in targeted in situ spatial transcriptomics mixes three things: what the cell intrinsically is, what varies with its microenvironment, and transcripts that were assigned to it but belong to its neighbours. The second and third are confounded at first order, since both make a cell's expression predictable from its neighbours' types, so no model fit to a single section can tell them apart without further assumptions. We present DISCELL, a variational model that splits a cell's expression into an intrinsic latent state $\vz$, a latent response $\vw$ associated with a deterministic description of its microenvironment, and a fixed leakage mixture over graph neighbours governed by a leak fraction $\kappa$. Type-conditional invariance keeps niche-predictable information out of $\vz$ through an adversarial regulariser, with a closed-form Gaussian penalty as the cheaper first attempt, and is graded by an independently trained held-out probe rather than by the training objective. Because $\kappa$ is not identifiable from a single section, we do not estimate it: every model is refitted on a grid of $\kappa$ values, and the deliverable is which spatial effects survive the grid. \todo{empirical claims to be added once the evaluation exists}
